% GDZ PPN236018647,native598–599;HKUST599–600.
\clearpage\SourceLeaf{598}\ArithmeticHead{569}{R}{1808 IUN. 12—1812 FEBR. 29.}
\NumberedItem{139}
\noindent\hspace*{\parindent}Series ad Media arithmetico geometrica pertinentes fusius evolutae.
\par\noindent\hfill 1809 Iun. 20.
\par\medskip\begingroup\small
\noindent\hspace*{\parindent}Man kann annehmen, daß der Brief \textsc{Schumachers} an \textsc{Gauss} vom 2. April 1808 (abgedruckt oben S. 242, [4.]), namentlich die darin erwähnte Aufgabe des \textsc{Pedrayes} (vergl. \textsc{Gauss}' Antwortschreiben vom 17. September 1808, abgedruckt oben S. 243, [5.]) die äußere Veranlassung dazu geboten hat, daß \textsc{Gauss} die Untersuchungen über elliptische Funktionen, die seit 1800 zurückgetreten waren, wieder aufnahm. An die hier zu besprechende Tagebuchaufzeichnung erinnert der Anfang der im III. Bande der Werke, S. 446 abgedruckten nachgelassenen Abhandlung, die S. 221—233 des im Oktober 1805 begonnenen Handbuchs 18, Bd aufgezeichnet ist. Diese Abhandlung beginnt nämlich mit den Worten: »Die Theoreme in Beziehung auf diejenigen Reihen und unendlichen Producte, welche zu der Theorie der Arithmetisch Geometrischen Mittel gehören, ordnen wir so«. Da sie unmittelbar auf eine astronomische Rechnung folgt, der die Bemerkung beigefügt ist: »geendiget d. 28. April 1809« (vergl. \textsc{Scherings} Bemerkung, Werke III, S. 494) ist ihr Zusammenhang mit unserer Tagebuchnotiz gesichert. Dasselbe gilt von der oben S 213 abgedruckten Aufzeichnung, die den Seiten 37—40 der Scheda An entnommen ist; auf S. 35 der Scheda, am Schluß einer astronomischen Rechnung findet sich nämlich die Angabe: »geendigt d. 2. May 1809«.
\par\noindent\hfill\textsc{Klein.\quad Schlesinger.}
\par\endgroup

\NumberedItem{140}
\noindent\hspace*{\parindent}Quinquesectionem pro mediis arithm[etico] Geom[etricis] absol[vimus].
\par\noindent\hfill 1809 Iun. 29.
\par\medskip\begingroup\small
\noindent\hspace*{\parindent}Die Fünfteilung der Perioden wird in der bei der Nr. 139 genannten Abhandlung des Handbuchs 18 Bd, siehe Werke III, S. 456 ff. entwickelt.
\par\noindent\hfill\textsc{Klein.\quad Schlesinger.}
\par\endgroup

\NumberedItem{141}
\noindent\hspace*{\parindent}Catalogum praecedentem per fata iniqua iterum interruptum initio anni 1812 resumimus. In mense Nov. 1811 contigerat demonstrationem theorematis fundamentalis in doctrina aequationum pure analyticam completam reddere; sed quum nihil chartis servatum fuerit, pars quaedam essentialis memoriae penitus exciderat. Hanc per satis longum temporis intervallum frustra quaesitam tandem feliciter redinvenimus.
\par\noindent\hfill 1812 Febr. 29.
\par\medskip\begingroup\small
\noindent\hspace*{\parindent}Es handelt sich um die \textit{Demonstratio nova altera} des Fundamentalsatzes der Algebra, vorgelegt am 7. Dezember 1815, Commentationes soc. reg. sc. Gottingensis rec. \textbf{3}, 1816, Werke III, S. 31; vergl. die Anzeige ebenda, S. 105. In einem Briefe an \textsc{Olbers} vom 19. Februar 1826, \textit{Wilhelm Olbers, Sein Leben und seine Werke} II2, Berlin 1909, S. 439, schreibt \textsc{Gauss}:
\par\endgroup
\par\medskip\noindent\textbf{X1.}\hfill72
\clearpage\SourceLeaf{599}\ArithmeticHead{570}{L}{TAGEBUCH MIT ERLÄUTERUNGEN.}
\noindent\hspace*{\parindent}Ich habe in meinem wissenschaftlichen Leben öfters den Fall gehabt, dass ich durch äussere Umstände veranlasst, Beschäftigungen, die nicht glückten, bei Seite legte, und die allerdings später glückten, z. B. mein Beweis für das Haupttheorem der Lehre von den Gleichungen, der in dem 3. Bande unsrer Comment[ationes] steht; aber ich habe nachher die 10fache Anstrengung gehabt, nur erst wieder auf den Punkt zu kommen, auf dem ich schon früher mehr als einmahl gewesen war.
\par\begingroup\small
\noindent\hspace*{\parindent}Vergl. auch die Bemerkung von M. \textsc{Brendel}, Werke VII, 1906, S. 610.
\par\noindent\hfill\textsc{Klein.\quad Schlesinger.}
\par\endgroup

\NumberedItem{142}
\noindent\hspace*{\parindent}Theoriam Attractionis Sphaeroidis Elliptici in puncta extra solidum sita prorsus novam invenimus.
\par\noindent\hfill Seeberg[ae], 1812 Sept. 26.
\par\medskip\begingroup\small
\noindent\hspace*{\parindent}Siehe die folgende Nummer. Die Ortsangabe bezieht sich auf die Sternwarte Seeberg bei Gotha.
\par\noindent\hfill\textsc{Klein.}
\par\endgroup

\NumberedItem{143}
\noindent\hspace*{\parindent}Etiam partes reliquas eiusdem theoriae per methodum novam mirae simplicitatis absolvimus.
\par\noindent\hfill 1812 Oct. 15.\quad Gott[ingae.]
\par\medskip\begingroup\small
\noindent\hspace*{\parindent}Die Abhandlung: \textit{Theoria attractionis corporum sphaeroidicorum ellipticorum homogeneorum, methodo nova tractata} wurde am 18. März 1813 der Königl. Societät d. W. übergeben und ist im zweiten Bande der Commentationes, Göttingen 1813 veröffentlicht worden; sie ist abgedruckt Werke V, S. 1—22. In der Selbstanzeige (Gött. gel. Anzeigen, 5. April 1813, Werke V, S. 281) heißt es: »Der Verfasser der gegenwärtigen Abhandlung, welcher seit lange schon die Überzeugung hatte, daß die \textls[120]{echte} Auflösungsmethode jener berühmten Aufgabe erst noch gefunden werden müsse, wurde vor einem halben Jahre veranlaßt, sich mit derselben näher zu beschäftigen«. Am 15. November 1812 schreibt \textsc{Gauss} an \textsc{Gerling} über diese Untersuchung das folgende:
\par\endgroup
\noindent\hspace*{\parindent}In der letzten Zeit habe ich mich mit der berühmten Aufgabe der Anziehung der elliptischen Sphäroide beschäftigt. Ich habe die Freude gehabt, eine Auflösung zu finden, deren Einfachheit und Eleganz alle meine Erwartungen noch übertroffen hat. Ich werde sie sobald es sich thun lässt in einer Vorlesung der Societät übergeben.
\par\begingroup\small
\noindent\hspace*{\parindent}Vergl. auch den Brief von \textsc{Gauss} an \textsc{Laplace} vom 5. November 1812, oben S. 379 und die zugehörigen Bemerkungen S. 380, sowie den Brief von \textsc{Gauss} an \textsc{Schumacher} vom 31. Dezember 1812, \textit{Briefwechsel} I, Altona 1860, S. 95.
\par\noindent\hfill\textsc{Klein.\quad Stäckel,}
\par\endgroup
